% main.tex -- AI-Guided Prediction of Thermoelastic Wave
% Propagation in Geological Media. Diploma defence deck, IITU 2026.

\documentclass[aspectratio=169,10pt]{beamer}
\usetheme{default}

\usefonttheme{professionalfonts}
\usepackage[T1]{fontenc}
\usepackage{lmodern}
\usepackage{amsmath}
\usepackage{amssymb}
\usepackage{booktabs}
\usepackage{array}
\usepackage{graphicx}

\definecolor{AcademicNavy}{HTML}{17324D}
\definecolor{AcademicBlue}{HTML}{2E5E88}
\definecolor{AcademicText}{HTML}{202124}
\definecolor{AcademicMuted}{HTML}{5F6B76}
\definecolor{AcademicRule}{HTML}{CBD5E1}
\definecolor{AcademicPanel}{HTML}{F6F8FA}

% Clean academic defence style: light background, dark text, restrained accents.
\setbeamertemplate{navigation symbols}{}
\setbeamertemplate{headline}{}
\setbeamertemplate{blocks}[default]
\setbeamersize{text margin left=8mm,text margin right=8mm}

\setbeamercolor{background canvas}{bg=white}
\setbeamercolor{normal text}{fg=AcademicText,bg=white}
\setbeamercolor{frametitle}{fg=AcademicNavy,bg=white}
\setbeamercolor{title}{fg=AcademicNavy,bg=white}
\setbeamercolor{subtitle}{fg=AcademicMuted,bg=white}
\setbeamercolor{author}{fg=AcademicText,bg=white}
\setbeamercolor{institute}{fg=AcademicMuted,bg=white}
\setbeamercolor{date}{fg=AcademicMuted,bg=white}
\setbeamercolor{structure}{fg=AcademicBlue}
\setbeamercolor{block title}{fg=AcademicNavy,bg=AcademicPanel}
\setbeamercolor{block body}{fg=AcademicText,bg=AcademicPanel}
\setbeamercolor{block title example}{fg=AcademicNavy,bg=AcademicPanel}
\setbeamercolor{block body example}{fg=AcademicText,bg=AcademicPanel}
\setbeamercolor{block title alerted}{fg=AcademicNavy,bg=AcademicPanel}
\setbeamercolor{block body alerted}{fg=AcademicText,bg=AcademicPanel}
\setbeamercolor{footline}{fg=AcademicMuted,bg=white}

\setbeamerfont{title}{series=\bfseries,size=\Large}
\setbeamerfont{subtitle}{size=\normalsize}
\setbeamerfont{frametitle}{series=\bfseries,size=\large}
\setbeamerfont{normal text}{size=\normalsize}
\setbeamerfont{block title}{series=\bfseries,size=\small}
\setbeamerfont{block body}{size=\small}
\setbeamerfont{footline}{size=\scriptsize}

\setbeamertemplate{itemize item}{\raise0.2ex\hbox{\tiny$\blacktriangleright$}}
\setbeamertemplate{itemize subitem}{\raise0.2ex\hbox{\tiny$\circ$}}
\setbeamertemplate{enumerate item}{\textbf{\insertenumlabel.}}

\setbeamertemplate{frametitle}{%
  \vspace{2mm}
  \insertframetitle\par
  \vspace{1mm}
  {\color{AcademicRule}\rule{\linewidth}{0.45pt}}
  \vspace{-1mm}
}

\setbeamertemplate{footline}{%
  \leavevmode
  \hbox{%
    \begin{beamercolorbox}[wd=.72\paperwidth,ht=2.5ex,dp=1ex,leftskip=8mm]{footline}%
      \color{AcademicMuted}\insertshorttitle
    \end{beamercolorbox}%
    \begin{beamercolorbox}[wd=.28\paperwidth,ht=2.5ex,dp=1ex,right,rightskip=8mm]{footline}%
      \color{AcademicMuted}\insertframenumber\,/\,\inserttotalframenumber
    \end{beamercolorbox}%
  }%
}

\setbeamertemplate{title page}{%
  \vspace{10mm}
  {\usebeamerfont{title}\usebeamercolor[fg]{title}\inserttitle\par}
  \vspace{3mm}
  {\color{AcademicRule}\rule{0.68\linewidth}{0.6pt}\par}
  \vspace{4mm}
  {\usebeamerfont{subtitle}\usebeamercolor[fg]{subtitle}\insertsubtitle\par}
  \vfill
  {\usebeamerfont{author}\usebeamercolor[fg]{author}\insertauthor\par}
  \vspace{2mm}
  {\usebeamerfont{institute}\usebeamercolor[fg]{institute}\insertinstitute\par}
  \vspace{2mm}
  {\usebeamerfont{date}\usebeamercolor[fg]{date}\insertdate\par}
}

% -- Metadata -----------------------------------
\title[AI-Guided Thermoelastic Prediction]{%
  AI-Guided Prediction of Thermoelastic Wave Propagation%
  \texorpdfstring{\\}{ }%
  in Geological Media%
}
\subtitle{A comparative Beamer presentation for diploma defence}

\author[Group of 4]{%
  Kasymkhanov Timur \and Askarov Insar \and Kaidrakhmetova Dinara \and Nurshanov Dias%
}
\institute{International Information Technology University}
\date{Bachelor's diploma defence, 2026}

\begin{document}

\begin{frame}[plain]
  \titlepage
  \vfill
  \begin{center}\footnotesize
    Research adviser: \textbf{Bakhyt N. Alipova} (PhD, Applied Mathematics)
  \end{center}
\end{frame}

\section{Motivation and Problem}

\begin{frame}{Relevance of the Topic}
  \begin{block}{Why this topic matters}
    \begin{itemize}
      \item Thermally affected rocks are relevant to geothermal systems,
            underground engineering, and geophysical interpretation.
      \item Temperature change may alter stress, deformation, and wave behavior
            inside the geological medium.
      \item Repeated coupled simulations are informative, but expensive for
            comparative directional studies.
    \end{itemize}
  \end{block}
  \vskip 6pt
  \begin{exampleblock}{Working position of this thesis}
    This thesis is positioned as a \emph{research prototype} for
    \emph{AI-Guided Prediction of Thermoelastic Wave Propagation in Geological Media}
    propagation. It compares neural surrogate approaches under simplified
    assumptions rather than replacing numerical solvers.
  \end{exampleblock}
\end{frame}

\begin{frame}{Research Problem}
  \begin{exampleblock}{Problem statement}
    The research problem is to estimate directional characteristics of
    thermoelastic wave propagation in geological media using an
    AI-Guided Prediction of Thermoelastic Wave Propagation in Geological Media
    under simplified physical assumptions.
  \end{exampleblock}
  \vskip 6pt
  \begin{itemize}
    \item The problem is important because coupled thermal and mechanical
          calculations can be costly when many directions or materials are compared.
    \item This work estimates selected source--probe responses and travel-time
          characteristics in a controlled model-based calculation setting.
    \item It does not claim a complete field-scale thermoelastic simulator or
          field-validated prediction system.
  \end{itemize}
\end{frame}

\begin{frame}{Aim and Tasks}
  \begin{block}{Aim}
    To develop and justify the
    \emph{AI-Guided Prediction of Thermoelastic Wave Propagation in Geological Media}
    framework.
  \end{block}
  \vskip 4pt
  \begin{enumerate}
    \item Establish the physical and geological basis of the problem.
    \item Formulate a simplified thermoelastic prediction task in two dimensions (2D).
    \item Build a COMSOL-derived synthetic dataset for controlled comparison.
    \item Compare Physics-Informed Neural Network (PINN), Fourier Neural
          Operator (FNO), MeshGraphNet, and Transformer-style surrogates.
    \item Interpret the results, limitations, and next development steps.
  \end{enumerate}
  \vskip 4pt
  \footnotesize
  The calculations were performed for representative geological materials,
  including granite, basalt, sandstone, and limestone.
\end{frame}

\begin{frame}{Object and Subject of the Study}
  \begin{columns}[T]
    \begin{column}{0.48\textwidth}
      \begin{block}{Object}
        Thermoelastic wave propagation in geological media represented by
        effective rock parameters.
      \end{block}
    \end{column}
    \begin{column}{0.48\textwidth}
      \begin{block}{Subject}
        AI-guided prediction of thermal and displacement
        response along a selected source--probe path.
      \end{block}
    \end{column}
  \end{columns}
  \vskip 6pt
  \begin{alertblock}{Scope}
    The work does not claim to solve the full thermoelastic problem in
    real geological formations. It studies a simplified two-dimensional (2D)
    \emph{comparative framework}.
  \end{alertblock}
\end{frame}

\begin{frame}{Thermoelastic Coupling Logic}
  \begin{columns}[T]
    \begin{column}{0.54\textwidth}
      \begin{enumerate}
        \item \textbf{Thermal excitation:} a local source changes the temperature field.
        \item \textbf{Thermal expansion:} temperature change produces local strain.
        \item \textbf{Stress formation:} constrained expansion creates thermal stress.
        \item \textbf{Mechanical response:} stress gradients drive displacement.
      \end{enumerate}
      \[
        T \longrightarrow \varepsilon_{\mathrm{th}}
        \longrightarrow \sigma \longrightarrow u
      \]
    \end{column}
    \begin{column}{0.42\textwidth}
      \footnotesize
      \begin{block}{Physical meaning}
        Temperature change causes thermal strain; thermal strain contributes
        to stress; stress affects displacement and wave propagation.
      \end{block}
      \scriptsize
      $T$ -- temperature [K]; $\varepsilon_{\mathrm{th}}$ -- thermal strain
      [dimensionless]; $\sigma$ -- stress [Pa]; $u$ -- displacement [m].
    \end{column}
  \end{columns}
\end{frame}

\section{Physical Basis}

\begin{frame}{Geological Media and Effective Parameters}
  \begin{columns}[T]
    \begin{column}{0.58\textwidth}
      \footnotesize
      \renewcommand{\arraystretch}{1.15}
      \begin{tabular}{l|r|r|r|r}
        \hline
        \textbf{Rock} & $\rho$ & $V_P$ & $k$ & $C_p$ \\
        & kg/m$^3$ & m/s & W/(m K) & J/(kg K) \\
        \hline
        Granite    & 2650 & 5850 & 2.50 & 850 \\
        Basalt     & 3000 & 6550 & 1.95 & 1300 \\
        Sandstone  & 2725 & 5000 & 2.25 & 1000 \\
        Limestone  & 2675 & 5400 & 2.55 & 1050 \\
        \hline
      \end{tabular}
      \vskip 6pt
      \scriptsize
      Effective scalar parameters support comparison under a locally
      homogeneous isotropic approximation, not field calibration.
    \end{column}
    \begin{column}{0.36\textwidth}
      \includegraphics[width=\textwidth,height=0.42\textheight,keepaspectratio]{img/sawn_limestone_texture}
      \vskip 2pt
      \scriptsize
      Public-domain limestone texture, Wikimedia Commons / Titus Tscharntke.
    \end{column}
  \end{columns}
\end{frame}

\begin{frame}{Mechanical Rock Parameters}
  \scriptsize
  \renewcommand{\arraystretch}{1.12}
  \setlength{\tabcolsep}{2.6pt}
  \begin{tabular}{l|l|r|r|r|r|r|r}
    \hline
    \textbf{Material} & \textbf{Type} & $\rho$ & $E$ & $\nu$ & $\mu$ & $K$ & $V_P/V_S$ \\
    & & kg/m$^3$ & GPa & -- & GPa & GPa & m/s \\
    \hline
    Granite & intrusive & 2650 & 76.28 & 0.245 & 30.63 & 49.84 & 5850/3400 \\
    Granodiorite & intrusive & 2795 & 85.91 & 0.255 & 34.24 & 58.35 & 6100/3500 \\
    Basalt & volcanic & 3000 & 103.96 & 0.266 & 41.07 & 73.95 & 6550/3700 \\
    Diabase & hypabyssal & 2975 & 98.22 & 0.274 & 38.56 & 72.36 & 6450/3600 \\
    Gabbro & mafic & 3075 & 114.27 & 0.287 & 44.40 & 89.33 & 6950/3800 \\
    Sandstone & sedimentary & 2725 & 55.75 & 0.259 & 22.13 & 38.61 & 5000/2850 \\
    Limestone & carbonate & 2675 & 61.48 & 0.277 & 24.08 & 45.90 & 5400/3000 \\
    Marble & metamorphic & 2725 & 82.41 & 0.270 & 32.43 & 59.82 & 6150/3450 \\
    Schist & foliated & 2800 & 68.20 & 0.267 & 26.91 & 48.82 & 5500/3100 \\
    Quartzite & quartz-rich & 2650 & 83.50 & 0.250 & 33.40 & 55.70 & 6150/3550 \\
    \hline
  \end{tabular}
  \vskip 6pt
  Elastic constants are effective midpoint values derived under isotropic
  assumptions for model-ready comparison.
\end{frame}

\begin{frame}{Thermal and Porosity Parameters}
  \scriptsize
  \renewcommand{\arraystretch}{1.12}
  \setlength{\tabcolsep}{3.2pt}
  \begin{tabular}{l|r|r|r|r|r|>{\raggedright\arraybackslash}p{0.22\textwidth}}
    \hline
    \textbf{Material} & $\alpha$ & $k$ & $C_p$ & $C$ & $\gamma$ & \textbf{Porosity} \\
    & $10^{-6}$/K & W/(m K) & J/(kg K) & MJ/(m$^3$ K) & MPa/K & \% \\
    \hline
    Granite & 7.9 & 2.50 & 850 & 2.25 & 1.18 & total 45.5 \\
    Granodiorite & -- & 2.40 & -- & -- & -- & -- \\
    Basalt & -- & 1.95 & 1300 & 3.90 & -- & total 19.0 \\
    Diabase & -- & 2.50 & 850 & 2.53 & -- & -- \\
    Gabbro & -- & 2.30 & 1000 & 3.08 & -- & total 43.5 \\
    Sandstone & 11.6 & 2.25 & 1000 & 2.73 & 1.34 & total 31.5; eff. 26.5 \\
    Limestone & 8.0 & 2.55 & 1050 & 2.81 & 1.10 & total 31.5; eff. 18.0 \\
    Marble & 9.8 & 2.80 & 850 & 2.32 & 1.76 & -- \\
    Schist & -- & 2.75 & 1150 & 3.22 & -- & total 26.5; eff. 27.5 \\
    Quartzite & -- & 5.10 & 1000 & 2.65 & -- & -- \\
    \hline
  \end{tabular}
  \vskip 6pt
  Missing thermal expansion or coupling values are shown explicitly as dashes
  rather than inferred.
\end{frame}

\begin{frame}{Rock Properties Controlling Propagation}
  \begin{columns}[T]
    \begin{column}{0.48\textwidth}
      \begin{block}{Mechanical controls}
        \begin{itemize}
          \item Density controls inertia.
          \item Elastic stiffness controls resistance to deformation.
          \item Porosity, cracks, and saturation change effective stiffness.
        \end{itemize}
      \end{block}
    \end{column}
    \begin{column}{0.48\textwidth}
      \begin{block}{Thermal controls}
        \begin{itemize}
          \item Conductivity controls heat transfer.
          \item Heat capacity controls thermal storage.
          \item Thermal expansion converts temperature change into strain.
        \end{itemize}
      \end{block}
    \end{column}
  \end{columns}
  \vskip 6pt
  Directional behavior can appear because bedding, foliation, aligned
  pores, fractures, and stress state make real rocks structurally ordered.
\end{frame}

\begin{frame}{Elastic Waves in Geological Media}
  \begin{columns}[T]
    \begin{column}{0.48\textwidth}
      \begin{block}{Compressional (P) waves}
        Compressional waves. Particle motion is mainly parallel to the
        propagation direction and involves volume change.
      \end{block}
    \end{column}
    \begin{column}{0.48\textwidth}
      \begin{block}{Shear (S) waves}
        Shear waves. Particle motion is mainly transverse to the direction
        of propagation and depends on shear rigidity.
      \end{block}
    \end{column}
  \end{columns}
  \vskip 6pt
  \[
    V_P = \sqrt{\frac{K+\frac{4}{3}G}{\rho}},
    \qquad
    V_S = \sqrt{\frac{G}{\rho}}
  \]
  \scriptsize
  $V_P,V_S$ -- compressional and shear wave velocities [m/s];
  $K$ -- bulk modulus [Pa]; $G$ -- shear modulus [Pa];
  $\rho$ -- density [kg/m$^3$]. These relations show that stiffness
  increases wave velocity, while density increases inertia.
\end{frame}

\begin{frame}{Heat Conduction and Thermal Expansion}
  \begin{block}{Fourier heat flux and heat equation}
    \[
      q = -k\nabla T,
      \qquad
      \rho C_p\frac{\partial T}{\partial t}
      =
      \nabla\cdot(k\nabla T)+Q
    \]
    \scriptsize
    $q$ -- heat flux [W/m$^2$]; $k$ -- thermal conductivity [W/(m K)];
    $T$ -- temperature [K]; $\rho$ -- density [kg/m$^3$];
    $C_p$ -- specific heat capacity [J/(kg K)]; $t$ -- time [s];
    $Q$ -- volumetric heat source [W/m$^3$].
  \end{block}
  \vskip 4pt
  \begin{block}{Temperature perturbation and thermal strain}
    \[
      \theta = T-T_0,
      \qquad
      \varepsilon_{\mathrm{th}}=\alpha\Delta T
    \]
    \scriptsize
    $\theta,\Delta T$ -- temperature change [K]; $T_0$ -- reference
    temperature [K]; $\varepsilon_{\mathrm{th}}$ -- thermal strain
    [dimensionless]; $\alpha$ -- linear thermal expansion [1/K].
  \end{block}
\end{frame}

\begin{frame}{Key Physical Parameters}
  \begin{block}{Thermal diffusivity relation}
    \[
      \alpha_T = \frac{k}{\rho C_p}
    \]
  \end{block}
  \vskip 6pt
  \begin{itemize}
    \item $\alpha_T$ is thermal diffusivity [m$^2$/s].
    \item $k$ is thermal conductivity [W/(m K)].
    \item $\rho$ is density [kg/m$^3$], and $C_p$ is heat capacity [J/(kg K)].
    \item Larger $\alpha_T$ means faster spreading of a temperature disturbance.
  \end{itemize}
\end{frame}

\begin{frame}{Thermoelastic Coupling}
  \begin{columns}[T]
    \begin{column}{0.58\textwidth}
      \begin{block}{Thermoelastic extension of Hooke's law}
        \[
          \theta = T - T_0
        \]
        \[
          \sigma_{ij} = \lambda\,\delta_{ij}\,\varepsilon_{kk}
          + 2\mu\,\varepsilon_{ij}
          - \gamma\,\delta_{ij}\,\theta,
          \qquad
          \gamma = 3K\alpha
        \]
      \end{block}
      \begin{block}{Equation of motion}
        \[
          \rho\,\frac{\partial^2u_i}{\partial t^2}
          = \frac{\partial \sigma_{ij}}{\partial x_j} + f_i
        \]
      \end{block}
    \end{column}
    \begin{column}{0.38\textwidth}
      \scriptsize
      $\theta$ -- temperature perturbation [K]; $T,T_0$ -- current and
      reference temperatures [K]; $\sigma_{ij}$ -- stress tensor [Pa];
      $\varepsilon_{ij}$ -- strain tensor [dimensionless];
      $\varepsilon_{kk}$ -- volumetric strain [dimensionless];
      $\lambda,\mu,K$ -- Lame and bulk moduli [Pa];
      $\delta_{ij}$ -- Kronecker delta [dimensionless];
      $\alpha$ -- thermal expansion [1/K]; $\gamma$ -- thermoelastic
      coupling coefficient [Pa/K]; $u_i$ -- displacement [m];
      $x_j$ -- coordinate [m]; $f_i$ -- body force density [N/m$^3$].
      \vskip 4pt
      The relation states that stress depends on elastic deformation and
      temperature-induced expansion; this is why thermal and elastic fields
      are coupled in the model.
    \end{column}
  \end{columns}
\end{frame}

\begin{frame}{Mathematical Scope and Primary Fields}
  \begin{columns}[T]
    \begin{column}{0.48\textwidth}
      \begin{block}{Simplified model scope}
        \footnotesize
        \begin{itemize}
          \item 2D continuum: $\Omega\subset\mathbb{R}^{2}$,
                $(x,t)\in\Omega\times[0,T_f]$.
          \item In-plane displacement: $u(x,t)=[u,v]^T$.
          \item Small deformation and linear isotropic thermoelasticity.
          \item Pores, fractures, fluids, plasticity, damage, and full
                anisotropy are not resolved explicitly.
        \end{itemize}
      \end{block}
    \end{column}
    \begin{column}{0.48\textwidth}
      \footnotesize
      \begin{tabular}{@{}lll@{}}
        \toprule
        \textbf{Field} & \textbf{Meaning} & \textbf{Role} \\
        \midrule
        $T(x,t)$ & Temperature & Thermal state \\
        $\theta(x,t)$ & Perturbation & Thermal loading \\
        $(u,v)$ & Displacement & Probe response \\
        $\varepsilon_{ij}$ & Strain & Deformation \\
        $\sigma_{ij}$ & Stress & Internal force \\
        $d,\hat d,\phi$ & Direction & Source--probe query \\
        \bottomrule
      \end{tabular}
      \vskip 6pt
      These fields connect the physical formulation with the prediction
      methodology used in the comparative prototype.
    \end{column}
  \end{columns}
  \vskip 4pt
  \scriptsize
  These assumptions keep the task suitable for comparative AI-assisted
  directional prediction and make the results interpretable within the scope
  of a bachelor thesis. Here $\Omega$ is the model domain [m$^2$], $x$ is
  position [m], $t$ is time [s], $T_f$ is final time [s], and $u,v$ are
  displacement components [m].
\end{frame}

\begin{frame}{Directional Prediction Task}
  \begin{columns}[T]
    \begin{column}{0.56\textwidth}
      \includegraphics[width=\textwidth]{img/web_visualization}
    \end{column}
    \begin{column}{0.40\textwidth}
      \scriptsize
      A source point and a probe point define the directional query.
      \[
        d = x_p - x_s,\qquad
        L = \|d\|,\qquad
        \hat d = \frac{d}{\|d\|}
      \]
      \[
        \phi = \operatorname{atan2}(d_y,d_x),
        \qquad
        \mathcal{F}:X\rightarrow\widehat{Y}
      \]
      \scriptsize
      $x_s,x_p$ -- source and probe positions [m]; $d$ -- direction vector
      [m]; $L$ -- distance [m]; $\hat d$ -- unit direction [dimensionless];
      $\phi$ -- azimuth [rad]; $\mathcal{F}$ maps input features $X$ to
      predicted outputs $\widehat{Y}$.
      \begin{block}{Prediction contract}
        \scriptsize
        Input: material parameters, geometry, direction, and time.
        Output: estimated temperature, displacement, directional response,
        and travel time.
      \end{block}
    \end{column}
  \end{columns}
\end{frame}

\section{Prediction Framework}

\begin{frame}{COMSOL-Based Synthetic Dataset}
  \begin{block}{Controlled thermoelastic calculation}
    Two-dimensional (2D) plane-strain rock domain $1\,\mathrm{m}\times1\,\mathrm{m}$ with a
    heated internal source used to generate time-dependent synthetic data.
  \end{block}
  \vskip 4pt
  \begin{itemize}
    \item Initial rock temperature: $293.15$ K
    \item Source temperature: $1500$ K
    \item Outputs exported per node and time step: $T$, $(u,v)$, velocity,
          stress, strain, and material features
    \item Purpose: create a common reference dataset for model comparison
  \end{itemize}
\end{frame}

\begin{frame}{Data Representation and Preprocessing}
  \begin{itemize}
    \item COMSOL comma-separated value (CSV) exports are aligned into a single node set.
    \item Dynamic features are assembled into a spatio-temporal tensor:
          \[
            \mathbf{S}\in\mathbb{R}^{N_t \times N_p \times F_d}
          \]
    \item Static material and geometry features are stored separately.
    \item Feature channels are normalized before training:
          \[
            \widetilde{x}=\frac{x-\mu}{\sigma+\varepsilon}
          \]
    \item The same processed dataset supports coordinate, grid, and graph models.
  \end{itemize}
  \scriptsize
  $\mathbf{S}$ is the assembled data tensor; $N_t$ is the number of time
  steps, $N_p$ the number of spatial points, and $F_d$ the number of dynamic
  feature channels. $\widetilde{x}$ is a normalized feature [dimensionless],
  $x$ is the original feature in its native unit, $\mu$ and $\sigma$ are its
  dataset mean and standard deviation, and $\varepsilon$ is a small
  dimensionless stabilizer.
\end{frame}

\section{Comparative Framework}

\begin{frame}{Comparative Prototype Framework}
  \begin{block}{One workflow for all model routes}
    The frontend, backend orchestrator, material catalog, and model
    services use one shared input--output contract for fair comparison.
  \end{block}
  \vskip 6pt
  \begin{itemize}
    \item Same source--probe geometry for all routes
    \item Same effective geological material presets
    \item Same normalized response format: temperature, displacement,
          direction, travel time, and diagnostics
  \end{itemize}
\end{frame}

\begin{frame}{Shared Model Comparison Contract}
  \begin{block}{Common interpretation}
    All routes are treated as surrogate modelling approaches for the same
    AI-Guided Prediction of Thermoelastic Wave Propagation in Geological Media
    task, not as complete numerical solvers.
  \end{block}
  \begin{columns}[T]
    \begin{column}{0.32\textwidth}
      \begin{block}{Input}
        \footnotesize
        Material parameters, source--probe geometry, observation time, and
        processed field features.
      \end{block}
    \end{column}
    \begin{column}{0.32\textwidth}
      \begin{block}{Output}
        \footnotesize
        Estimated temperature, displacement components, travel time, and
        diagnostic response quantities.
      \end{block}
    \end{column}
    \begin{column}{0.32\textwidth}
      \begin{block}{Comparison basis}
        \footnotesize
        Same material presets, geometry contract, response format, and
        analytical travel-time sanity check.
      \end{block}
    \end{column}
  \end{columns}
\end{frame}

\begin{frame}{Physics-Informed Neural Network (PINN)}
  \begin{columns}[T]
    \begin{column}{0.42\textwidth}
      \footnotesize
      \begin{block}{Role}
        Physics-informed baseline in the comparison.
      \end{block}
      \begin{block}{Input / output}
        Input: coordinates, time, and material parameters.
        Output: estimated temperature and in-plane displacement.
      \end{block}
      \begin{block}{Why and limitation}
        Used because thermoelastic residuals can be included in training.
        It is still a baseline model, not an exact equation solver.
      \end{block}
    \end{column}
    \begin{column}{0.54\textwidth}
      \begin{block}{Model relation}
        \[
          f_{\theta}(x,y,t,E,\nu,\rho,\alpha,k,C_p)\rightarrow(T,u,v)
        \]
        \[
          \mathcal{L}
          =
          \lambda_{\mathrm{sup}}\mathcal{L}_{\mathrm{sup}}
          +
          \lambda_{\mathrm{vel}}\mathcal{L}_{\mathrm{vel}}
          +
          \lambda_{\mathrm{el}}\mathcal{L}_{\mathrm{elastic}}
          +
          \lambda_{\mathrm{th}}\mathcal{L}_{\mathrm{thermal}}
        \]
      \end{block}
      \scriptsize
      $x,y$ -- coordinates [m]; $t$ -- time [s]; $E$ -- Young's modulus [Pa];
      $\nu$ -- Poisson's ratio [dimensionless]; $\rho$ -- density [kg/m$^3$];
      $\alpha$ -- thermal expansion [1/K]; $k$ -- conductivity [W/(m K)];
      $C_p$ -- heat capacity [J/(kg K)]; $T,u,v$ -- predicted temperature [K]
      and displacements [m]. $\mathcal{L}$ and $\mathcal{L}_{\cdot}$ are
      normalized loss terms; $\lambda_{\cdot}$ are dimensionless loss weights.
    \end{column}
  \end{columns}
\end{frame}

\begin{frame}{Fourier Neural Operator (FNO)}
  \begin{columns}[T]
    \begin{column}{0.48\textwidth}
      \begin{block}{Role}
        Neural operator surrogate for structured two-dimensional fields.
      \end{block}
      \begin{block}{Input}
        \footnotesize
        Grid-based material and response features, together with the shared
        source--probe and time query context.
      \end{block}
    \end{column}
    \begin{column}{0.48\textwidth}
      \begin{block}{Output}
        \footnotesize
        Estimated field-style response used to derive directional quantities.
      \end{block}
      \begin{block}{Why and limitation}
        \footnotesize
        Included because Fourier operators can capture broad field patterns.
        In the current prototype, its output scale needs recalibration.
      \end{block}
    \end{column}
  \end{columns}
\end{frame}

\begin{frame}{FNO: Field-to-Field Prediction Idea}
  \begin{columns}[T]
    \begin{column}{0.48\textwidth}
      \begin{block}{Why FNO fits this task}
        \footnotesize
        Thermoelastic simulation output is not a single scalar. It is a set
        of spatial fields defined on a two-dimensional grid.
      \end{block}
      \begin{itemize}
        \footnotesize
        \item The model learns an operator between function spaces.
        \item Inputs and outputs keep spatial structure.
        \item The target can include temperature and displacement fields.
      \end{itemize}
    \end{column}
    \begin{column}{0.48\textwidth}
      \begin{block}{Operator view}
        \[
          \mathcal{G}_{\theta}:\mathcal{A}\rightarrow\mathcal{U},
          \qquad
          \hat{s}=\mathcal{G}_{\theta}(a)
        \]
      \end{block}
      \scriptsize
      $a$ -- input field; $\hat{s}$ -- predicted thermoelastic field;
      $\mathcal{A}$ and $\mathcal{U}$ -- input and output field spaces;
      $\theta$ -- trainable neural-operator parameters.
    \end{column}
  \end{columns}
\end{frame}

\begin{frame}{FNO Tensor Format for the 2D Setup}
  \begin{columns}[T]
    \begin{column}{0.50\textwidth}
      \begin{block}{Input tensor}
        \[
          A\in\mathbb{R}^{B\times C_{in}\times N_x\times N_y}
        \]
      \end{block}
      \begin{block}{Output tensor}
        \[
          \hat{S}\in\mathbb{R}^{B\times C_{out}\times N_x\times N_y}
        \]
      \end{block}
    \end{column}
    \begin{column}{0.46\textwidth}
      \footnotesize
      \begin{itemize}
        \item $B$ is batch size.
        \item $N_x,N_y$ are grid sizes.
        \item $C_{in}$ is the number of input channels.
        \item $C_{out}$ is controlled by \texttt{out\_channels}.
      \end{itemize}
      \vskip 4pt
      For the three-channel thermoelastic setup:
      \[
        C_{out}=3,\qquad \hat{s}=[\hat{T},\hat{u},\hat{v}].
      \]
    \end{column}
  \end{columns}
\end{frame}

\begin{frame}{FNO Architecture Used in the Prototype}
  \begin{block}{Main processing chain}
    \[
      \text{input}\rightarrow\text{lifting}\rightarrow
      \text{Fourier blocks}\rightarrow\text{projection}\rightarrow\text{output}
    \]
  \end{block}
  \vskip 4pt
  \begin{columns}[T]
    \begin{column}{0.47\textwidth}
      \footnotesize
      \begin{block}{Implementation parameters}
        \begin{itemize}
          \item \texttt{modes}: retained Fourier modes
          \item \texttt{num\_fourier\_layers}: number of blocks
          \item \texttt{mid\_channels}: latent capacity
          \item \texttt{projection\_channels}: output projection size
        \end{itemize}
      \end{block}
    \end{column}
    \begin{column}{0.49\textwidth}
      \begin{block}{Operator composition}
        \[
          \mathcal{G}_{\theta}=Q\circ\mathcal{B}_{L-1}\circ\cdots
          \circ\mathcal{B}_0\circ P
        \]
      \end{block}
      \scriptsize
      $P$ -- lifting layer; $\mathcal{B}_l$ -- Fourier block;
      $Q$ -- projection layer; $L$ -- number of Fourier blocks.
    \end{column}
  \end{columns}
\end{frame}

\begin{frame}{FNO Spectral Block}
  \begin{columns}[T]
    \begin{column}{0.50\textwidth}
      \begin{block}{Block update}
        \[
          v_{l+1}=\sigma\left(\mathcal{S}_l(v_l)+\mathcal{C}_l(v_l)\right)
        \]
      \end{block}
      \footnotesize
      The spectral branch captures global frequency-domain interactions;
      the local convolution branch captures neighboring spatial patterns.
    \end{column}
    \begin{column}{0.46\textwidth}
      \begin{block}{Spectral convolution}
        \[
          \hat{v}_l=\mathcal{F}(v_l),\qquad
          z_l=\mathcal{F}^{-1}\left(R_l\cdot\hat{v}_l\right)
        \]
      \end{block}
      \scriptsize
      $\mathcal{F}$ and $\mathcal{F}^{-1}$ are Fourier and inverse Fourier
      transforms; $R_l$ is the trainable complex-valued spectral weight tensor.
    \end{column}
  \end{columns}
\end{frame}

\begin{frame}{Mode Truncation and Coordinate Encoding}
  \begin{columns}[T]
    \begin{column}{0.48\textwidth}
      \begin{block}{Retained modes}
        \[
          \texttt{modes}=[K_1,K_2],\qquad d=2
        \]
      \end{block}
      \footnotesize
      Only selected low-frequency Fourier coefficients are transformed by
      learned weights. This controls capacity and computational cost.
    \end{column}
    \begin{column}{0.48\textwidth}
      \begin{block}{Optional grid channels}
        \[
          \tilde{a}(x,y)=[a(x,y),x,y]
        \]
      \end{block}
      \footnotesize
      With \texttt{add\_grid}, the model receives explicit normalized spatial
      coordinates. This helps distinguish source, boundary, and interior regions.
    \end{column}
  \end{columns}
\end{frame}

\begin{frame}{Training, Evaluation, and Current FNO Caveat}
  \begin{columns}[T]
    \begin{column}{0.52\textwidth}
      \begin{block}{Supervised objective}
        \[
          \mathcal{L}=\lambda_{MSE}\mathcal{L}_{MSE}
          +\lambda_{rel}\mathcal{L}_{rel}
        \]
      \end{block}
      \scriptsize
      Channel normalization is computed on the training set and reused for
      validation, testing, and inference.
      \[
        \tilde{s}_{i,c}=\frac{s_{i,c}-\mu_c}{\sigma_c}
      \]
    \end{column}
    \begin{column}{0.44\textwidth}
      \begin{alertblock}{Interpretation in this thesis}
        \footnotesize
        FNO is valuable as a field-based neural-operator route, but the current
        prototype still needs output-scale recalibration before it can be read
        as a physically reliable predictor.
      \end{alertblock}
      \footnotesize
      Evaluation should use field-level plots, channel-wise errors, relative
      $L^2$ error, and travel-time sanity checks.
    \end{column}
  \end{columns}
\end{frame}

\begin{frame}{MeshGraphNet}
  \begin{columns}[T]
    \begin{column}{0.48\textwidth}
      \begin{block}{Role}
        Graph-based surrogate model for mesh-aware prediction.
      \end{block}
      \begin{block}{Input}
        \footnotesize
        Mesh nodes, graph connectivity, material features, and query context.
      \end{block}
    \end{column}
    \begin{column}{0.48\textwidth}
      \begin{block}{Output}
        \footnotesize
        Node-level or probe-level response estimates after graph message passing.
      \end{block}
      \begin{block}{Why and limitation}
        \footnotesize
        Included because graph structure can represent local mesh interactions.
        The current route behaves mainly as a fixed-horizon predictor.
      \end{block}
    \end{column}
  \end{columns}
\end{frame}

\begin{frame}{Transformer-Style Surrogate Model}
  \begin{columns}[T]
    \begin{column}{0.48\textwidth}
      \begin{block}{Role}
        Attention-based baseline surrogate for fast query prediction.
      \end{block}
      \begin{block}{Input}
        \footnotesize
        Tokens describing geometry, material parameters, state variables,
        and requested observation time.
      \end{block}
    \end{column}
    \begin{column}{0.48\textwidth}
      \begin{block}{Output}
        \footnotesize
        Fast estimates of temperature, displacement, travel time, and diagnostics.
      \end{block}
      \begin{block}{Why and limitation}
        \footnotesize
        Included as a flexible low-latency route. Its physics priors are limited,
        so results require cautious interpretation.
      \end{block}
    \end{column}
  \end{columns}
\end{frame}

\section{Calculation Results}

\begin{frame}{Calculation Status}
  \begin{block}{Calculation setup}
    \begin{enumerate}
      \item \textbf{40 scenarios per material}, four materials:
            granite, sandstone, limestone, and basalt.
      \item \textbf{$\mathbf{40\,\times\,4\,\times\,4 = 640}$ model
            responses} evaluated under one shared contract.
      \item \textbf{Zero fallback}: every plotted response is an active
            model output, not a mock.
      \item \textbf{2D plane-strain only}: the operating envelope of
            all trained surrogates.
    \end{enumerate}
  \end{block}
  \begin{alertblock}{FNO caveat}
    \footnotesize
    FNO is included in the comparison, but its output scales are still
    miscalibrated. It should currently be treated as a diagnostic branch.
  \end{alertblock}
\end{frame}

\begin{frame}{Travel-Time Predictions}
  \begin{columns}[T]
    \begin{column}{0.62\textwidth}
      \includegraphics[width=\textwidth]{img/travel_time_by_model_material}
    \end{column}
    \begin{column}{0.34\textwidth}
      \footnotesize
      \begin{block}{Reading the plot}
        \begin{itemize}
          \item PINN $\approx 0.11$\,ms
          \item Transformer $\approx 0.14$\,ms
          \item FNO $\approx 1.29$\,ms
          \item MeshGraphNet $\approx 2.47$\,ms
        \end{itemize}
      \end{block}
      Model-family differences are larger than the basalt--sandstone
      material difference in the current prototype.
    \end{column}
  \end{columns}
\end{frame}

\begin{frame}{Accuracy Against Analytical Reference}
  \begin{columns}[T]
    \begin{column}{0.56\textwidth}
      \includegraphics[width=\textwidth]{img/travel_time_relative_error}
    \end{column}
    \begin{column}{0.40\textwidth}
      \footnotesize
      \begin{block}{Analytical reference}
      \[
        t_{\mathrm{gt}} = \frac{r}{V_p}
      \]
      \scriptsize
      $t_{\mathrm{gt}}$ -- reference travel time [s]; $r$ -- source--probe
      distance [m]; $V_p$ -- compressional wave velocity [m/s].
      \end{block}
      \scriptsize
      \begin{tabular}{@{}lcc@{}}
        \toprule
        \textbf{Model} & \textbf{Median error} & \textbf{$\leq$10\,\%} \\
        \midrule
        PINN         & $2.0\%$    & $100\%$ \\
        Transformer  & $29.2\%$   & $0\%$   \\
        FNO          & $1403.6\%$ & $0\%$   \\
        MeshGraphNet & $> 100\%$  & $0\%$   \\
        \bottomrule
      \end{tabular}
      \vskip 4pt
      PINN is the only route that remains inside the $10\%$ band across
      the balanced sweep.
    \end{column}
  \end{columns}
\end{frame}

\begin{frame}{Inference Latency and Time Behavior}
  \begin{columns}[T]
    \begin{column}{0.62\textwidth}
      \includegraphics[width=\textwidth]{img/inference_latency_by_model}
    \end{column}
    \begin{column}{0.34\textwidth}
      \footnotesize
      \begin{block}{Median latency}
        \begin{itemize}
          \item Transformer $\approx 1.3$\,ms
          \item MeshGraphNet $\approx 1.5$\,ms
          \item FNO $\approx 2.4$\,ms
          \item PINN $\approx 2.6$\,ms
        \end{itemize}
      \end{block}
      \vskip 4pt
      PINN is also the only route that changes smoothly with requested
      observation time; the other routes currently behave as fixed-horizon
      predictors.
    \end{column}
  \end{columns}
\end{frame}

\begin{frame}{Threats to Validity}
  \begin{enumerate}
    \item \textbf{Limited training data.}
          The COMSOL-derived pilot configuration is not a full calibrated
          geological database.
    \item \textbf{Simplified 2D assumptions.}
          The study uses effective material parameters in a plane-strain setup.
    \item \textbf{FNO normalization issue.}
          The current 2D adaptation is not yet physically reliable.
    \item \textbf{Analytical reference is isotropic.}
          $r/V_p$ is a useful sanity check, not a complete thermoelastic validation.
  \end{enumerate}
\end{frame}

\section{Conclusion}

\begin{frame}{Conclusions}
  \begin{enumerate}
    \item The thesis develops a physically grounded
          \emph{comparative framework} for
          AI-Guided Prediction of Thermoelastic Wave Propagation in Geological Media.
    \item PINN is the strongest \emph{physics-informed baseline} in the
          current comparison and shows the best agreement with the analytical
          travel-time reference.
    \item Transformer is the fastest route, while MeshGraphNet remains a
          useful mesh-aware surrogate candidate.
    \item FNO should currently be treated as a diagnostic branch until
          output normalization is corrected.
  \end{enumerate}
\end{frame}

\begin{frame}{Future Work}
  \begin{enumerate}
    \item Retrain and recalibrate the 2D FNO route.
    \item Add time-conditioned rollout behavior to MeshGraphNet.
    \item Extend the material catalog with more complete thermoelastic
          parameters and validation data.
    \item Compare against additional analytical or numerical baselines.
  \end{enumerate}
  \vskip 8pt
  \begin{exampleblock}{Final message}
    The prototype estimates directional characteristics of thermoelastic
    wave propagation within the simplified assumptions of the study.
  \end{exampleblock}
\end{frame}

\begin{frame}{Thank You}
  \centering
  \Large Questions?
\end{frame}

\end{document}
